% Created 2026-10-17 Sat 14:12
% Intended LaTeX compiler: pdflatex
\documentclass[11pt,letterpaper]{article}
\usepackage[utf8]{inputenc}
\usepackage[T1]{fontenc}
\usepackage[spanish]{babel}
\usepackage{graphicx}
\usepackage{amsmath}
\usepackage{amssymb}
\usepackage{hyperref}
\newcommand{\materia}{MODELADO Y PROGRAMACIÓN}
\newcommand{\tipoPractica}{DOCUMENTACIÓN}
\newcommand{\practicaTitulo}{PRÁCTICA 4:\\PRUEBAS UNITARIAS}
\newcommand{\practicaSubtitulo}{DISEÑO DE PRUEBAS PARA UN SISTEMA}
\newcommand{\autorUno}{Alan Gael Alvarez Hernández}
\newcommand{\autorDos}{María Fernanda Miranda Sánchez}
\newcommand{\autorTres}{Pedro Pablo Ruiz Hernández}
\newcommand{\nombreProfesora}{Karla Adriana Esquivel Guzmán}
\newcommand{\ayudanteTeoria}{Ossiel Aguilar Espíndola}
\newcommand{\ayudanteLaboratorio}{Enoc Pineda Aguirre}
\newcommand{\semestreFinal}{SEMESTRE 1-2027}
\newcommand{\fechaPortada}{Ciudad de México, 20 de octubre de 2026}
\newcommand{\leyendaEscudoLineaUno}{FACULTAD DE CIENCIAS}
\newcommand{\leyendaEscudoLineaDos}{UNAM}
\newcommand{\encabezadoIzquierdo}{Modelado y Programación}
\newcommand{\encabezadoDerecho}{Práctica 5: TDD}
\newif\ifportada
\portadatrue
\newif\ifbibliografia
\bibliografiatrue
\usepackage[T1]{fontenc}
\usepackage[utf8]{inputenc}
\usepackage{lmodern}
\usepackage{microtype}
\usepackage[normalem]{ulem}
\newcommand{\fuentePortada}{\fontfamily{ptm}\selectfont}
\newcommand{\fuenteInstitucional}{\fontfamily{ptm}\fontseries{m}\selectfont}
\AtBeginDocument{\renewcommand{\contentsname}{Índice}\renewcommand{\tablename}{Tabla}\renewcommand{\figurename}{Figura}\renewcommand{\refname}{Bibliografía}}
\AtBeginDocument{\decimalpoint}
\AtBeginDocument{\shorthandoff{"}}
\usepackage[top=2.2cm,bottom=2.2cm,left=2.3cm,right=2.3cm]{geometry}
\setlength{\parindent}{0cm}
\setlength{\parskip}{0.65em}
\emergencystretch=2em
\usepackage{graphicx}
\usepackage{float}
\usepackage{caption}
\usepackage{array}
\usepackage{longtable}
\usepackage{ragged2e}
\usepackage{makecell}
\newcolumntype{L}[1]{>{\RaggedRight\arraybackslash\sloppy\hspace{0pt}}p{#1}}
\newcolumntype{C}[1]{>{\Centering\arraybackslash}p{#1}}
\usepackage{enumitem}
\setlist[itemize,1]{label=\textbullet}
\setlist[itemize,2]{label=--}
\setlist[itemize,3]{label=$\circ$}
\setlist[enumerate,1]{label=\arabic*.}
\setlist[enumerate,2]{label=\alph*)}
\usepackage{xcolor}
\usepackage{tikz}
\usetikzlibrary{calc}
\usepackage{listings}
\lstdefinestyle{practica}{
basicstyle=\ttfamily\small,
backgroundcolor=\color{gray!7},
frame=single,
rulecolor=\color{black},
framerule=0.45pt,
framesep=5pt,
xleftmargin=0pt,
xrightmargin=0pt,
aboveskip=0.4em,
belowskip=0.8em,
showstringspaces=false,
columns=fullflexible,
keepspaces=true,
breaklines=true,
literate={á}{{\'a}}1 {é}{{\'e}}1 {í}{{\'i}}1 {ó}{{\'o}}1 {ú}{{\'u}}1 {Á}{{\'A}}1 {É}{{\'E}}1 {Í}{{\'I}}1 {Ó}{{\'O}}1 {Ú}{{\'U}}1 {ñ}{{\~n}}1 {Ñ}{{\~N}}1
}
\usepackage{titlesec}
\setcounter{secnumdepth}{2}
\setcounter{tocdepth}{2}
\titleformat{\section}{\Large\bfseries}{}{0pt}{}
\titleformat{\subsection}{\large\bfseries}{}{0pt}{}
\titlespacing*{\section}{0pt}{1.5em}{0.8em}
\titlespacing*{\subsection}{0pt}{1.2em}{0.5em}
\usepackage{fancyhdr}
\setlength{\headheight}{15pt}
\pagestyle{fancy}
\fancyhf{}
\fancyhead[L]{\encabezadoIzquierdo}
\fancyhead[R]{\encabezadoDerecho}
\fancyfoot[C]{\thepage}
\renewcommand{\headrulewidth}{0.4pt}
\renewcommand{\footrulewidth}{0pt}
\hypersetup{colorlinks=true,linkcolor=black,urlcolor=blue,citecolor=black}
\author{Nombre del alumno}
\date{\today}
\title{Reporte}
\hypersetup{
 pdfauthor={Nombre del alumno},
 pdftitle={Reporte},
 pdfkeywords={},
 pdfsubject={},
 pdfcreator={Emacs 30.1 (Org mode 9.7.11)}, 
 pdflang={Spanish}}
\begin{document}


\ifportada

\hypersetup{pageanchor=false}

\begin{titlepage}

\thispagestyle{empty}


\begin{tikzpicture}[remember picture,overlay]


% ============================================================
% ESCUDO UNAM
% ============================================================

\IfFileExists{imagenes/unam.png}{

    \node[
        anchor=north west,
        inner sep=0pt
    ]
    at ([xshift=2.92cm,yshift=-2.17cm]current page.north west)
    {
        \includegraphics[width=3.68cm]{imagenes/unam.png}
    };

}{}


% ============================================================
% UNIVERSIDAD NACIONAL AUTÓNOMA DE MÉXICO
%
% La referencia usa Times con iniciales manualmente más grandes.
% Iniciales: ~16 pt.
% Resto: ~13 pt.
% ============================================================

\node[
    anchor=north,
    align=center,
    font=\fuenteInstitucional
]
at ([xshift=12.10cm,yshift=-2.51cm]current page.north west)
{
    {\fontsize{16.02}{18.4}\selectfont U}{\fontsize{13.02}{18.4}\selectfont NIVERSIDAD }
    {\fontsize{16.02}{18.4}\selectfont N}{\fontsize{13.02}{18.4}\selectfont ACIONAL }
    {\fontsize{16.02}{18.4}\selectfont A}{\fontsize{13.02}{18.4}\selectfont UTÓNOMA}\\[0.12cm]
    {\fontsize{13.02}{18.4}\selectfont DE }
    {\fontsize{16.02}{18.4}\selectfont M}{\fontsize{13.02}{18.4}\selectfont ÉXICO}
};


% ============================================================
% DOBLE LÍNEA HORIZONTAL
% ============================================================

% Línea principal

\draw[line width=1.90pt]
([xshift=7.00cm,yshift=-3.98cm]current page.north west)
--
([xshift=17.05cm,yshift=-3.98cm]current page.north west);


% Línea secundaria

\draw[line width=0.60pt]
([xshift=7.00cm,yshift=-4.50cm]current page.north west)
--
([xshift=17.05cm,yshift=-4.50cm]current page.north west);


% ============================================================
% FACULTAD DE CIENCIAS
%
% Misma construcción tipográfica de la referencia:
% iniciales grandes + resto ligeramente menor.
% ============================================================

\node[
    anchor=north,
    align=center,
    font=\fuenteInstitucional
]
at ([xshift=12.10cm,yshift=-5.31cm]current page.north west)
{
    {\fontsize{16.02}{18.0}\selectfont F}{\fontsize{13.02}{18.0}\selectfont ACULTAD DE }
    {\fontsize{16.02}{18.0}\selectfont C}{\fontsize{13.02}{18.0}\selectfont IENCIAS}
};


% ============================================================
% TRES LÍNEAS VERTICALES
% ============================================================

% Línea izquierda

\draw[line width=0.60pt]
([xshift=4.22cm,yshift=-6.27cm]current page.north west)
--
([xshift=4.22cm,yshift=-21.35cm]current page.north west);


% Línea central

\draw[line width=1.95pt]
([xshift=4.78cm,yshift=-6.27cm]current page.north west)
--
([xshift=4.78cm,yshift=-21.35cm]current page.north west);


% Línea derecha

\draw[line width=0.60pt]
([xshift=5.35cm,yshift=-6.27cm]current page.north west)
--
([xshift=5.35cm,yshift=-21.35cm]current page.north west);


% ============================================================
% MATERIA
% ============================================================

\node[
    anchor=north,
    align=center,
    text width=11.5cm,
    font=\fuentePortada\fontsize{13.0}{15.2}\selectfont
]
at ([xshift=12.10cm,yshift=-7.78cm]current page.north west)
{
    \materia
};


% ============================================================
% TIPO DE PRÁCTICA
% ============================================================

\node[
    anchor=north,
    align=center,
    font=\fuentePortada\fontsize{10.8}{12.8}\selectfont
]
at ([xshift=12.10cm,yshift=-8.60cm]current page.north west)
{
    \tipoPractica
};


% ============================================================
% TÍTULO DE LA PRÁCTICA
% ============================================================

\node[
    anchor=north,
    align=center,
    text width=11.8cm,
    font=\fuentePortada\fontsize{20.8}{24.2}\selectfont
]
at ([xshift=12.10cm,yshift=-10.36cm]current page.north west)
{
    \practicaTitulo
};


% ============================================================
% SUBTÍTULO DE LA PRÁCTICA
% ============================================================

\node[
    anchor=north,
    align=center,
    text width=11.6cm,
    font=\fuentePortada\fontsize{10.8}{12.8}\selectfont
]
at ([xshift=12.10cm,yshift=-12.42cm]current page.north west)
{
    \practicaSubtitulo
};


% ============================================================
% PRESENTAN
% ============================================================

\node[
    anchor=north,
    align=center,
    font=\fuentePortada\fontsize{11.3}{13.4}\selectfont
]
at ([xshift=12.10cm,yshift=-14.42cm]current page.north west)
{
    P\hspace{0.31cm}
    R\hspace{0.31cm}
    E\hspace{0.31cm}
    S\hspace{0.31cm}
    E\hspace{0.31cm}
    N\hspace{0.31cm}
    T\hspace{0.31cm}
    A\hspace{0.31cm}
    N
    \hspace{0.31cm}:
};


% ============================================================
% EQUIPO
%
% Alineado con las líneas horizontales superiores:
% ancho = 10.05 cm.
%
% Separaciones:
% encabezado -> primer integrante = 0.52 cm
% entre integrantes = 0.44 cm
% último integrante -> PROFESORA = 1.04 cm
%
% Los datos se editan al inicio del archivo:
% \autorUno / \cuentaUno
% \autorDos / \cuentaDos
% \autorTres / \cuentaTres
% ============================================================

% Encabezados
% Mismo tamaño y altura tipográfica que PROFESORA:
\node[
    anchor=north,
    align=center,
    font=\fuentePortada\mdseries\fontsize{10.8}{12.8}\selectfont
]
at ([xshift=12.025cm,yshift=-16.25cm]current page.north west)
{
    INTEGRANTES DEL EQUIPO:
};


% Integrante 1
% La separación visual EQUIPO -> primer integrante es la misma
% que PROFESORA -> nombre de profesora.
\node[
    anchor=north,
    align=center,
    text width=12.2cm,
    font=\fuentePortada\mdseries\fontsize{9.8}{11.6}\selectfont
]
at ([xshift=12.025cm,yshift=-16.9cm]current page.north west)
{
    \autorUno
};


% Integrante 2
\node[
    anchor=north,
    align=center,
    text width=12.2cm,
    font=\fuentePortada\mdseries\fontsize{9.8}{11.6}\selectfont
]
at ([xshift=12.025cm,yshift=-17.34cm]current page.north west)
{
    \autorDos
};


% Integrante 3
\node[
    anchor=north,
    align=center,
    text width=12.2cm,
    font=\fuentePortada\mdseries\fontsize{9.8}{11.6}\selectfont
]
at ([xshift=12.025cm,yshift=-17.78cm]current page.north west)
{
    \autorTres
};

% ============================================================
% PROFESORA
% ============================================================

\node[
    anchor=north,
    align=center,
    font=\fuentePortada\fontsize{10.8}{12.8}\selectfont
]
at ([xshift=12.10cm,yshift=-18.82cm]current page.north west)
{
    PROFESORA:
};

\node[
    anchor=north,
    align=center,
    text width=12.2cm,
    font=\fuentePortada\mdseries\fontsize{9.8}{11.6}\selectfont
]
at ([xshift=12.10cm,yshift=-19.32cm]current page.north west)
{
    \nombreProfesora
};


% ============================================================
% AYUDANTE DE TEORÍA
% ============================================================

\node[
    anchor=north,
    align=center,
    font=\fuentePortada\fontsize{10.0}{11.9}\selectfont
]
at ([xshift=12.10cm,yshift=-20.36cm]current page.north west)
{
    AYUDANTE DE TEORÍA:
};

\node[
    anchor=north,
    align=center,
    text width=12.2cm,
    font=\fuentePortada\mdseries\fontsize{9.8}{11.6}\selectfont
]
at ([xshift=12.10cm,yshift=-20.88cm]current page.north west)
{
    \ayudanteTeoria
};


% ============================================================
% AYUDANTE DE LABORATORIO
% ============================================================

\node[
    anchor=north,
    align=center,
    font=\fuentePortada\fontsize{10.0}{11.9}\selectfont
]
at ([xshift=12.10cm,yshift=-21.92cm]current page.north west)
{
    AYUDANTE DE LABORATORIO:
};

\node[
    anchor=north,
    align=center,
    text width=12.2cm,
    font=\fuentePortada\mdseries\fontsize{9.8}{11.6}\selectfont
]
at ([xshift=12.10cm,yshift=-22.44cm]current page.north west)
{
    \ayudanteLaboratorio
};


% ============================================================
% ESCUDO FACULTAD DE CIENCIAS
% ============================================================

\IfFileExists{imagenes/ciencias.png}{

    \node[
        anchor=north west,
        inner sep=0pt
    ]
    at ([xshift=3.49cm,yshift=-22.20cm]current page.north west)
    {
        \includegraphics[width=2.65cm]{imagenes/ciencias.png}
    };

    % ========================================================
    % LEYENDA DEBAJO DEL ESCUDO DE CIENCIAS
    %
    % Se deja en dos comandos separados para que sea fácil cambiar,
    % por ejemplo, a otra facultad.
    % ========================================================

    % Primera línea de la leyenda:
    % centrada con el escudo y muy cerca de la imagen.
    \node[
        anchor=north,
        align=center,
        font=\fuentePortada\fontsize{4.05}{4.60}\selectfont
    ]
    at ([xshift=4.815cm,yshift=-25.12cm]current page.north west)
    {
        \leyendaEscudoLineaUno
    };

    % Segunda línea:
    % separada de forma independiente para que UNAM quede
    % pegado a FACULTAD DE CIENCIAS como en la referencia.
    \node[
        anchor=north,
        align=center,
        font=\fuentePortada\fontsize{3.70}{4.20}\selectfont
    ]
    at ([xshift=4.815cm,yshift=-25.33cm]current page.north west)
    {
        \leyendaEscudoLineaDos
    };

}{}


% ============================================================
% SEMESTRE
% ============================================================

\node[
    anchor=north,
    align=center,
    font=\fuentePortada\mdseries\fontsize{11.0}{13.0}\selectfont
]
at ([xshift=12.10cm,yshift=-24.12cm]current page.north west)
{
    \semestreFinal
};


% ============================================================
% LUGAR Y FECHA
%
% Se deja editable al inicio del archivo con \fechaPortada.
% ============================================================

\node[
    anchor=north,
    align=center,
    font=\fuentePortada\mdseries\fontsize{10.4}{12.4}\selectfont
]
at ([xshift=12.10cm,yshift=-24.78cm]current page.north west)
{
    \fechaPortada
};


\end{tikzpicture}


\end{titlepage}


% ============================================================
% REINICIO DE NUMERACIÓN
% ============================================================

\setcounter{page}{2}

\hypersetup{pageanchor=true}

\fi


% ============================================================
% NOMBRES EN ESPAÑOL
% ============================================================

\renewcommand{\contentsname}{Índice}

\renewcommand{\tablename}{Tabla}

\renewcommand{\figurename}{Figura}

\renewcommand{\refname}{Bibliografía}


% ============================================================
% ÍNDICE
% ============================================================

\tableofcontents

\clearpage
\section{Introducción}
\label{sec:org3b51c88}

\subsection{Pruebas Unitarias}
\label{sec:org3063f07}

Una \textbf{Prueba Unitaria} verifica un comportamiento pequeño del sistema: una función, un método o una regla concreta. Para el desarrollo de las mismas se recomienda seguir el siguiente flujo:


\begin{figure}[H]

\centering

{\large
\textbf{Problema}
\,$\longrightarrow$\,
\textbf{Modelo}
\,$\longrightarrow$\,
\textbf{Casos}
\,$\longrightarrow$\,
\textbf{Pruebas}
\,$\longleftrightarrow$\,
\textbf{Implementación}
}

\caption{Flujo general del diseño de pruebas.}

\end{figure}

\textbf{Nota:} El resultado esperado se decide antes de ejecutar el código. No se copia lo que devuelve la implementación.

Los casos de prueba definidos se clasifican como:

\begin{itemize}
\item \textbf{Normal:} Un valor dentro de la regla.
\item \textbf{Frontera:} Justo antes, en y después de un límite.
\item \textbf{Inválido:} Un valor que el sistema debe rechazar.
\item \textbf{Excepción:} Un error esperado.
\item \textbf{Cruce de reglas:} Cuando dos condiciones se cumplen al mismo tiempo.
\end{itemize}

\textbf{Ejemplo:}

Si hasta 15 minutos es gratis, y desde 16 se cobra, probar 14, 15, 16 y 17.
\subsection{Estructura de una prueba Arrange - Act - Asseert}
\label{sec:org1e57857}

\begin{itemize}
\item \textbf{Arrange:} Preparación de los datos, objetos, o estado inicial.
\item \textbf{Act}: Ejecución de la función o método que se desea probar.
\item \textbf{Assert:} Comprobación de que lo obtenido sea lo esperado.
\end{itemize}

\textbf{Ejemplo:}


\begin{lstlisting}[style=practica]
# Arrange
minutos = 16

# Act
resultado = sistema.calcular_total(minutos, "normal", False)

# Assert
assert resultado == 20.0
\end{lstlisting}

\textbf{Nota:} Una buena prueba responde tres preguntas: \textbf{¿qué regla comprueba?, ¿qué entrada utiliza? y ¿qué resultado debería producir?}
\section{Objetivo}
\label{sec:orgd4851fc}

El alumno diseñará un conjunto de pruebas unitarias para verificar el comportamiento de un sistema a partir de su especificación, no únicamente a partir de lo que hace el código.

Al finalizar la práctica, el alumno deberá ser capaz de:

\begin{itemize}
\item identificar casos normales;
\item identificar casos frontera;
\item detectar entradas inválidas;
\item verificar errores y excepciones;
\item escribir pruebas usando el patrón Arrange–Act–Assert;
\item justificar de qué regla del modelo proviene cada caso de prueba;
\item ejecutar una suite de pruebas y analizar sus resultados;
\item diferenciar entre “la prueba pasó” y “el programa es correcto”.
\end{itemize}
\section{Desarrollo de la práctica}
\label{sec:orgdcb2f5e}

\subsection{Durante la práctica}
\label{sec:org05a9391}

\begin{enumerate}
\item Lea primero el modelo.
\item Diseñe los casos sin ejecutar todavía el sistema.
\item Escriba el resultado esperado de cada caso.
\item Implemente las pruebas.
\item Ejecute la suite.
\item Si alguna prueba falla, no cambie inmediatamente la prueba.
\item Compare:
\end{enumerate}

\begin{center}


\textbf{modelo esperado} vs. \textbf{resultado obtenido}
\end{center}

\begin{enumerate}
\item Determine si el problema esta en:
\begin{enumerate}
\item la implementación;
\item la prueba;
\item la interpretación del modelo.
\end{enumerate}
\end{enumerate}
\subsection{Plantemaiento del Problema}
\label{sec:orge387a45}

Probar el sistema de Ticket de estacionamiento.

El sistema recibe:

\begin{itemize}
\item minutos de permanencia;
\item tipo de cliente;
\item si perdió el boleto;
\end{itemize}
\subsubsection{Modelo del Sistema: Ticket de estacionamiento}
\label{sec:orgf85885d}

Comportamiento esperado:

\textbf{R1.} Los minutos deben ser un número entero mayor o igual que 0.

\textbf{R2.} Los tipos de clientes válidos son:

\begin{itemize}
\item \textbf{"normal"}
\item \textbf{"frecuente"}
\end{itemize}

\textbf{R3.} Si el boleto está perdido, el costo final es de \$300.00, sin importar los minutos ni el tipo de cliente.

\textbf{R4.} Si los minutos son 0, el costo es \$0.00

\textbf{R5.} De 1 a 15 minutos, el costo es \$0.00

\textbf{R6.} De 16 a 60 minutos, el costo es de \$20.00

\textbf{R7.} Después de 60 minutos se cobran \$20.00 por la primera hora y \$15.00 por cada hora adicional \textbf{iniciada}.

\textbf{R8.} Un cliente \textbf{"frecuente"} recibe un 10\% de descuento sobre el costo calculado, excepto cuando perdió el boleto.

\textbf{R9.} El costo nunca puede ser negativo.

\textbf{R10.} Si los minutos son negativos, el sistema debe lanzar un error.

\textbf{R11.} Si el tipo de cliente no es válido, el sistema debe lanzar un error.
\subsubsection{Ejemplos del Sistema: Ticket de estacionamiento}
\label{sec:orgdbf9ccc}

Ejemplos de interpretación:

\begin{itemize}
\item 15 minutos, \textbf{cliente normal} → \$0.00
\item 16 minutos, \textbf{cliente normal} → \$20.00
\item 60 minutos, \textbf{cliente normal} → \$20.00
\item 61 minutos, \textbf{cliente normal} → \$35.00
\item 120 minutos, \textbf{cliente normal} → \$35.00
\item 121 minutos, \textbf{cliente normal} → \$50.00
\item 61 minutos, \textbf{cliente frecuente} → \$20.00
\item \emph{boleto perdido} → \$300.00
\end{itemize}
\subsection{Planteamiento de Pruebas}
\label{sec:orga091683}

\subsection{Implementación en Python}
\label{sec:orga70e9c2}
\section{Checklist}
\label{sec:org65f5aad}

Ejemplo:

\begin{itemize}
\item[{$\boxtimes$}] Actividad terminada
\item[{$\boxminus$}] Actividad parcialmente realizada
\item[{$\square$}] Una actividad
\end{itemize}

Actividades:

\begin{itemize}
\item[{$\boxminus$}] Documentación.
\item[{$\square$}] Planteamiento de Pruebas.
\item[{$\square$}] Implementación de Pruebas.
\item[{$\square$}] Evidencia de Ejecución.
\item[{$\square$}] Preguntas de Análisis.
\item[{$\square$}] Organización de Archivos.
\item[{$\square$}] README.md.
\end{itemize}
\end{document}
